%%%%%%%%%%%%%%%%%%%%%%%%%%%%%%%%%%%%%%%%%%%%%%%%%%%%%%%%%%%%%%%%%%%%%%%%%%%%%%%%
%  IE643 course project poster -- A1 portrait (59.4 cm x 84.1 cm), two columns
%
%  Left column : Abstract, Introduction, Datasets, Methodology (pipeline diagram)
%  Right column: Novelty Assessment, Results (intensive vs. novelty),
%                Conclusion, References, QR codes
%
%  Compile with pdfLaTeX or XeLaTeX (e.g. upload the whole final_ie643_report
%  folder to Overleaf). Figures are read from figures/:
%     logo.png, IE643-Logo.pdf      header logos
%     poster_dataflow.png           pipeline diagram (scripts/make_poster_dataflow.py)
%     poster_generalisation.png     MetroPT-3 result plot
%     github_qr.png, demo_qr.png    your QR codes (add them yourself)
%  A missing figure is drawn as a grey box with its file name.
%
%  FITTING THE PAGE: each column is stretched to the page height (\vfill).
%  If a column runs off the bottom, lower \flowscale (left) or \plotscale
%  (right); if space is left over, raise them (max 1).
%
%  Numbers: experiments/benchmark_v2/summary.csv (base training) and
%  experiments/benchmark_v2_regularized/summary.csv (regularized training);
%  GRU rows are the mean of 3 runs.
%%%%%%%%%%%%%%%%%%%%%%%%%%%%%%%%%%%%%%%%%%%%%%%%%%%%%%%%%%%%%%%%%%%%%%%%%%%%%%%%
\documentclass[a1,portrait]{a0poster}
\usepackage[paperwidth=59.4cm, paperheight=84.1cm, margin=0cm]{geometry}
\usepackage[T1]{fontenc}
\usepackage[sfdefault]{roboto}
\usepackage[svgnames]{xcolor}
\usepackage{graphicx}
\usepackage{tikz}
\usepackage{booktabs}
\usepackage{colortbl}
\usepackage{adjustbox}
\usepackage{amsmath, amssymb}
\usepackage{enumitem}
\usepackage[hidelinks]{hyperref}
\graphicspath{{figures/}}
\urlstyle{same}                                          % links in the poster font
\expandafter\def\expandafter\UrlBreaks\expandafter{\UrlBreaks\do\-}   % long links may wrap at hyphens

%% ---- the two size knobs (fraction of the column width) ----
\newcommand{\flowscale}{1.00}   % pipeline diagram, left column
\newcommand{\plotscale}{0.62}   % MetroPT-3 plot, right column

\definecolor{HeadBlue}{HTML}{1F3864}
\definecolor{Ours}{HTML}{FDECEC}
\definecolor{Block}{HTML}{EEF2F7}
\setlength{\parindent}{0pt}
\setlength{\parskip}{0pt}
\setlength{\fboxsep}{0pt}
\renewcommand{\arraystretch}{1.15}
\setlist[itemize]{leftmargin=1.0em, itemsep=3pt, topsep=3pt, parsep=0pt,
                  label={\color{HeadBlue}\scriptsize$\blacksquare$}}
\setlist[enumerate]{leftmargin=1.4em, itemsep=3pt, topsep=3pt, parsep=0pt,
                    label={\color{HeadBlue}\bfseries\arabic*.}}
\newlength{\bodywidth}
\newlength{\colwidth}
\newlength{\colheight}

%% section heading
\newcommand{\sect}[1]{\par\vspace{0.1cm}{\huge\bfseries\color{HeadBlue}#1\par}\vspace{0.15cm}}
%% figure, or a grey box with the file name if the file is missing
\newcommand{\fig}[3]{\IfFileExists{figures/#1}{\includegraphics[width=#2]{#1}}%
  {\fcolorbox{Gray}{WhiteSmoke}{\parbox[c][#3][c]{#2}{\centering\color{DimGray}%
  {\tiny add\\ \texttt{\detokenize{#1}}}}}}}
\newcommand{\logobox}[1]{\tikz[baseline=(current bounding box.center)]\node[fill=white, rounded corners=0.45cm, inner sep=0.35cm]
  {\includegraphics[height=5.6cm, width=5.6cm, keepaspectratio]{#1}};}
%% small grey caption text
\newcommand{\figcap}[1]{\par\vspace{0.12cm}{\footnotesize\color{DimGray}#1\par}}
%% shaded block heading inside a table
\newcommand{\blockrow}[1]{\rowcolor{Block}\multicolumn{5}{@{}l}{\textit{#1}}\\}

\begin{document}
\pagestyle{empty}
%% page-dependent sizes (set here, once the A1 page size is in force)
\setlength{\bodywidth}{\dimexpr\paperwidth-3cm\relax}
\setlength{\colwidth}{\dimexpr(\bodywidth-2.2cm)/2\relax}
\setlength{\colheight}{\dimexpr\paperheight-12.4cm\relax}
%% everything sits in one page-sized box, so header and columns stay together
\noindent\begin{minipage}[t][\dimexpr\paperheight-0.4cm\relax][t]{\paperwidth}
%==============================================================================
%  HEADER
%==============================================================================
\noindent\colorbox{white}{\begin{minipage}[c][8.8cm][c]{\paperwidth}
\hspace*{1.6cm}\logobox{logo.png}\hfill
\begin{minipage}[c]{0.66\paperwidth}\centering\color{HeadBlue}
  {\normalsize IE 643 Course Project\par}\vspace{0.15cm}
  {\LARGE\bfseries Vision-Language Model Guided Knowledge Distillation\\[0.1cm]
   for Time-Series Anomaly Detection\par}\vspace{0.2cm}
  {\normalsize Team Ballad\par}\vspace{0.05cm}
  {\normalsize\color{Black} Nisith Naman Parida (24B2218)\par}\vspace{0.05cm}
  {\small\color{DimGray} 24b2218@iitb.ac.in\par}
\end{minipage}\hfill
\logobox{IE643-Logo.pdf}\hspace*{1.6cm}
\end{minipage}}\par
\noindent{\color{HeadBlue}\rule{\paperwidth}{6pt}}\par
\vspace{0.3cm}

%==============================================================================
%  BODY: two columns separated by a blue rule
%==============================================================================
\noindent\hspace*{1.5cm}%
%------------------------------------------------------------------ LEFT COLUMN
\begin{minipage}[t][\colheight][s]{\colwidth}
\small
\sect{Abstract}
Anomaly detection in multivariate time-series data is challenging due to complex temporal patterns and highly
imbalanced class distributions, where anomalous instances are significantly underrepresented. A pretrained
Vision-Language Model serves as a frozen teacher, providing anomaly-related guidance from visual representations
of time-series data. This knowledge is transferred to a trainable GRU-based student model through knowledge
distillation, enabling it to learn temporal dependencies from numerical sequences. The approach aims to achieve
effective anomaly detection and temporal localization with reduced computational requirements during inference.

\vfill
\sect{Introduction}
\begin{itemize}
  \item \textbf{Goal}: find \emph{whether} and \emph{where} sensors behave unlike normal (fraud, freezing of
        gait, oil-well faults, compressor failures).
  \item \textbf{Hard}: faults are rare (5--7\% of windows in three datasets) and costly to label.
  \item \textbf{Why CLIP}: many faults are visible shapes; CLIP knows irregular pictures from 400 million
        image--caption pairs, but has 151 million parameters.
  \item \textbf{Idea}: frozen CLIP teaches a small GRU (95,000 parameters) in training; only the GRU is deployed.
  \item \textbf{Question}: does the teacher help \emph{beyond} the same GRU trained on labels alone?
\end{itemize}

\vfill
\sect{Datasets}
{\footnotesize
\begin{adjustbox}{max width=\linewidth}
\begin{tabular}{@{}l l c c l c@{}}
\toprule
\textbf{Dataset} & \textbf{Measures} & \textbf{Sensors} & \textbf{Window} & \textbf{Test kept apart by} & \textbf{Faulty} \\
\midrule
Credit Card  & card payments         & 29 & 64  & time (60/20/20\%)  & 7.4\%  \\
Daphnet FoG  & walking, 64 Hz        & 9  & 256 & other patients     & 4.8\%  \\
Petrobras 3W & oil wells, 10 s       & 8  & 120 & other well events  & 62.0\% \\
MetroPT-3    & air compressor, 1 min & 15 & 120 & other failures     & 6.2\%  \\
\bottomrule
\end{tabular}
\end{adjustbox}\par}
\figcap{Window = time steps per piece; Faulty = share of faulty test windows. MetroPT-3: failure 1 trains,
2 validates, 3--4 test. Air Quality was dropped: its labels proved artificial.}

\vfill
\sect{Methodology}
{\centering\fig{poster_dataflow.png}{\flowscale\linewidth}{34cm}\par}
\figcap{Each card shows what the data looks like at that step (toy: 4 sensors, 24-step windows). Orange:
frozen CLIP teacher, training only; red: the deployed GRU student.}
\end{minipage}%
\hfill{\color{HeadBlue}\vrule width 5pt}\hfill
%------------------------------------------------------------------ RIGHT COLUMN
\begin{minipage}[t][\colheight][s]{\colwidth}
\small
\sect{Novelty Assessment}
\begin{enumerate}
  \item \textbf{Label audit}: Air Quality's faults were synthetic, so it was dropped; 3W and MetroPT-3 added.
  \item \textbf{Teacher input}: fixed $\pm4$ colour scale (v1 pictures were nearly flat) + a spectrogram.
  \item \textbf{Learned prompts}: 14 prompt words learned for frozen CLIP, cross-fitted (no self-scoring).
  \item \textbf{Confidence-weighted distillation} of CLIP's opinion and its picture summary.
  \item \textbf{Controlled benchmark}: 5 teachers $\times$ 7 students $\times$ 3 runs + 2 baselines.
  \item \textbf{Fair baseline}: regularized training so the labels-only GRU cannot memorise.
  \item \textbf{Checks}: metrics recomputed, leak and overfitting audits, unlabelled test with 12 hidden
        failures.
\end{enumerate}

\vfill
\sect{Results}
{\footnotesize
\begin{adjustbox}{max width=\linewidth}
\begin{tabular}{@{}l c c c c@{}}
\toprule
\textbf{Test window AUPRC} & \textbf{Credit Card} & \textbf{Daphnet} & \textbf{3W} & \textbf{MetroPT-3} \\
\midrule
\blockrow{Primary assessment (v1): saved runs, global colour range, hand-written prompts}
\multicolumn{5}{@{}l}{Daphnet: CLIP teacher 0.029 $\cdot$ GRU, labels only 0.126 $\cdot$ GRU + KD 0.058} \\
\multicolumn{5}{@{}l}{Credit Card: GRU + KD 0.660 $\cdot$ 3W and MetroPT-3 were not evaluated in v1} \\
\blockrow{Intensive assessment: original pipeline re-run on the final splits}
Random guess (= share of faults)        & 0.074 & 0.048 & 0.620 & 0.062 \\
Isolation Forest                        & 0.128 & 0.057 & 0.787 & 0.575 \\
Gradient-boosted trees                  & 0.767 & 0.350 & 0.998 & 0.702 \\
CLIP, hand-written prompts (teacher)    & 0.079 & 0.123 & 0.771 & 0.037 \\
GRU, labels only                        & \textbf{0.823} & \textbf{0.627} & \textbf{0.985} & \textbf{0.751} \\
GRU + KD from that teacher              & 0.817 & 0.609 & 0.982 & 0.424 \\
\blockrow{Novelty assessment: adapted teacher and new distillation}
CLIP, learned prompts                   & 0.137 & 0.531 & 0.974 & 0.989 \\
CLIP, learned prompts + spectrogram (teacher) & 0.124 & 0.401 & 0.971 & 0.998 \\
GRU + KD, learned-prompt teacher        & 0.804 & 0.587 & 0.978 & \textbf{0.893} \\
\quad + confidence weighting            & 0.801 & 0.592 & 0.978 & 0.876 \\
GRU, labels only, regularized           & \textbf{0.828} & \textbf{0.659} & \textbf{0.985} & 0.812 \\
\rowcolor{Ours}\textbf{GRU + KD, confidence-weighted, regularized (ours)} & 0.805 & 0.655 & 0.983 & 0.892 \\
\bottomrule
\end{tabular}
\end{adjustbox}\par}
\figcap{Table 1. Test window AUPRC (0--1, higher is better; random guessing scores the share of faults).
v1: single saved runs on the v1 splits. Other GRU rows: mean of 3 runs; bold = best GRU in each block.}
\vspace{0.3cm}
{\footnotesize
\begin{adjustbox}{max width=\linewidth}
\begin{tabular}{@{}l cc cc cc@{}}
\toprule
 & \multicolumn{2}{c}{\textbf{Window AUPRC}} & \multicolumn{2}{c}{\textbf{Time-step AUPRC}} & \multicolumn{2}{c}{\textbf{Event F1}} \\
\textbf{Regularized: labels only vs.\ ours} & labels & ours & labels & ours & labels & ours \\
\midrule
Credit Card & 0.828 & 0.805 & 0.777 & 0.776 & 0.798 & 0.792 \\
Daphnet     & 0.659 & 0.655 & 0.535 & 0.549 & 0.290 & 0.232 \\
3W          & 0.985 & 0.983 & 0.979 & 0.986 & 0.945 & 0.949 \\
\rowcolor{Ours}MetroPT-3 & 0.812 & \textbf{0.892} & 0.821 & \textbf{0.882} & 0.577 & \textbf{0.820} \\
\bottomrule
\end{tabular}
\end{adjustbox}\par}
\figcap{Table 2. Same network, split and seeds, with and without the teacher; no point adjustment.}
\begin{itemize}
  \item \textbf{The teacher had to be adapted}: the v1 teacher scored 0.029 on Daphnet and hand-written prompts
        0.037 on MetroPT-3 (distilling them hurts: 0.424); learned prompts reach 0.531 and 0.989.
  \item \textbf{MetroPT-3, one training failure}: the adapted teacher beats labels only on all 3 runs
        (+0.080 window, +0.061 time-step AUPRC, event F1 0.577 $\rightarrow$ 0.820).
  \item \textbf{New unlabelled data}, 12 hidden failures: both GRUs find 10; ours places them better
        (event F1 0.546 vs.\ 0.424).
  \item \textbf{Elsewhere no change} (within 0.025). Caveat: only two MetroPT-3 test failures.
\end{itemize}

\vfill
\sect{Conclusion and Future Work}
\begin{itemize}
  \item A frozen picture-and-text model \textbf{can} teach a sensor detector once its pictures and prompts are
        adapted; it helps \textbf{where labelled faults are scarcest}.
  \item The deployed GRU is \textbf{1/1500 of CLIP's size} and \textbf{330$\times$ faster} on a CPU.
  \item \textbf{Future}: leave-one-failure/patient-out tests; distil on unlabelled data; one prompt per fault type.
\end{itemize}

\vfill
\sect{References}
{\scriptsize
\begin{enumerate}[label={[\arabic*]}, leftmargin=2.0em, itemsep=2pt, topsep=0pt]
\item A.\ Radford et al. Learning transferable visual models from natural language supervision. ICML, 2021.
\item K.\ Zhou, J.\ Yang, C.\ C.\ Loy, Z.\ Liu. Learning to prompt for vision-language models. IJCV, 2022.
\item G.\ Hinton, O.\ Vinyals, J.\ Dean. Distilling the knowledge in a neural network. NIPS Workshop, 2015.
\item A.\ Romero et al. FitNets: hints for thin deep nets. ICLR, 2015.
\item S.\ Kim et al. Towards a rigorous evaluation of time-series anomaly detection. AAAI, 2022.
\item F.\ T.\ Liu, K.\ M.\ Ting, Z.-H.\ Zhou. Isolation Forest. IEEE ICDM, 2008.
\item M.\ B\"achlin et al. Wearable assistant for Parkinson's disease patients with the freezing of gait symptom. IEEE T-ITB, 2010.
\item R.\ E.\ V.\ Vargas et al. A realistic and public dataset with rare undesirable real events in oil wells. J.\ Pet.\ Sci.\ Eng., 2019.
\item B.\ Veloso, R.\ P.\ Ribeiro, J.\ Gama, P.\ M.\ Pereira. The MetroPT dataset for predictive maintenance. Scientific Data, 2022.
\item ULB Machine Learning Group. Credit Card Fraud Detection dataset. Kaggle.
\end{enumerate}
}
\vspace{0.25cm}
\begin{minipage}[c]{0.16\linewidth}\fig{github_qr.png}{\linewidth}{\linewidth}\end{minipage}\hspace{0.5cm}%
\begin{minipage}[c]{0.29\linewidth}\textbf{Code:}\\ GitHub repository\end{minipage}\hfill
\begin{minipage}[c]{0.16\linewidth}\fig{demo_qr.png}{\linewidth}{\linewidth}\end{minipage}\hspace{0.5cm}%
\begin{minipage}[c]{0.29\linewidth}\textbf{Demo:}\\ video of the app\end{minipage}
\par\vspace{0.12cm}
{\scriptsize\raggedright
\href{https://github.com/Nisith-Naman/Vision-Language-Model-Guided-Knowledge-Distillation-for-Time-Series-Anomaly-Detection}{\nolinkurl{github.com/Nisith-Naman/Vision-Language-Model-Guided-Knowledge-Distillation-for-Time-Series-Anomaly-Detection}}\par
\href{https://drive.google.com/file/d/1IPWDaN2lwlS4eK7C8tyMCtFCTHwT89Ao/view?usp=sharing}{\nolinkurl{drive.google.com/file/d/1IPWDaN2lwlS4eK7C8tyMCtFCTHwT89Ao/view?usp=sharing}}\par}
\end{minipage}\hspace*{1.5cm}
\end{minipage}
\end{document}
